\clearpage
\thispagestyle{empty}
\pdfbookmark[0]{Fundamentalgleichungen für die Bewegung schwerer Körper auf der rotierenden Erde. Ausgabe und Quellen}{rotating-earth-edition}
\begin{center}\Large\textbf{Fundamentalgleichungen für die Bewegung schwerer Körper auf der rotierenden Erde}\end{center}

Die beiden Briefe von Carl Friedrich Gauß an Johann Friedrich Benzenberg vom 2. Februar und 8. März 1803 und die vollständige Herleitung sind KI-Übertragungen des deutschen Textes der \textit{Werke}, Band V (Göttingen, 1867), S.~495--503, in modernes Deutsch. Der Satz einschließlich der bearbeitbaren Zeichnung wurde mit KI erstellt. Die Gleichungen und die angegebenen Zahlenwerte folgen der Werkausgabe.

Die Briefe und die Herleitung erschienen auch in Johann Friedrich Benzenbergs \textit{Dr. Benzenbergs Versuche über die Umdrehung der Erde} (Dortmund, bey Mallinkrodt, 1804), S.~349--353 bzw. 363--371. Das hier verlinkte Exemplar befindet sich in der ETH-Bibliothek Zürich (Rar 1913). Eine eigenständige moderne Darstellung von Fallversuchen geben S. Draszow und U. Backhaus in \textit{Dreht sich die Erde wirklich?} (Augsburg, 2003).

\medskip\noindent
\href{https://gdz.sub.uni-goettingen.de/id/PPN236006339}{Deutsche Werkausgabe (1867)}\par
\noindent\href{https://doi.org/10.3931/e-rara-14813}{Benzenbergs Buch (1804)}\par
\noindent\href{https://doi.org/10.1007/978-3-642-49319-5_23}{Deutscher Nachdruck (1877)}\par
\noindent\href{https://www.didaktik.physik.uni-due.de/backhaus/publicat/FallversucheDPG.pdf}{Draszow und Backhaus (2003)}
\clearpage
